\chapter{\label{app-reparationsopfer}Das Reparationsopfer}
tl;dr:
Das Reinigungsopfer wird für Sakrileg an Sankta dargebracht.
Andere Fälle sind auf diesen zurückzuführen.
Einer dieser Fälle ist der folgende: eine Person enteignet eine andere durch Diebstahl oder ähnliches.
Zu diesem zivilen Vergehen schwört der Täter einen falschen Schwur und belügt somit Gott.
Wenn die Person Reuhe empfindet und die Sünde bekennt, dann kann sie (a) Reparation an die enteignete Person machen und
(b) ein Reparationsopfer bringen, um den Sakrileg des falschen Schwurs auszutilgen.\footnote{
Buße verringert eine bewusste Lüge und falschen Schwur zu einem Vergehen, das wie ein unwissentlich Begangenes entfernt werden kann.
}
Danach vergibt Gott das Vergehen.
Hier ist ein erstaunlicher Fall, wo Buße und eine Reparationszahlung, sowie ein Reparationsopfer als Lösegeld für den Schuldigen Menschen fungieren.
In diesem Fall wird tatäschlich eine Sünde byd rmh durch Buße und ein Opfer vergeben.
Auch in diesem Fall wird keine Schuld auf das Opfer übertragen - es fungiert als Lösegeld, und wird nicht stellvertretend für den Opfernden bestraft.
